
\section{Bilinear Approximation v1 (BLA~v1)}
\label{sec:bilinear-approx}
\label{sec:approx-accel}

Perturbation theory (\cref{sec:perturbation-concept}) reduces the
precision required for per-pixel iteration, but it does not reduce the
\emph{number} of iterations: a pixel that does not escape still executes
the recurrence up to the requested limit~$n_{\max}$. At deep zooms that
limit can reach billions, and the work is repeated for every pixel.
The stored reference orbit has its own length~$n_{\mathrm{ref}}$;
rebasing permits evaluation to continue beyond that length. Neither
quantity is necessarily the orbit's period~$p$.

Approximation methods attack this remaining bottleneck by \emph{skipping}
large blocks of iterations entirely.  Instead of stepping through the
perturbation recurrence one iteration at a time, they precompute linear
maps that advance the perturbation delta by many iterations in a single
evaluation. \FractalShark{} implements two strategies:
\begin{itemize}
\item \textbf{Bilinear Approximation v1 (BLA~v1)}
  (\cref{sec:bilinear-approx})---precomputes composable two-term linear
  maps from the reference orbit, skipping iterations where the
  perturbations satisfy the step's validity test.
\item \textbf{Linear Approximation v2 (LA~v2)}
  (\cref{sec:la-v2-perturb})---precomputes staged linear-approximation
  coefficients that can skip entire sections of the orbit in one
  evaluation, achieving larger skips at the cost of more complex
  coefficient management.
\end{itemize}
Both techniques preserve the perturbation framework's mixed-precision
advantage while replacing many of the $n_{\max}$ individual updates
with block evaluations. The achievable reduction depends on the orbit,
the pixel offset, and the validity thresholds.

BLA~v1 is the original bilinear approximation path retained in
\FractalShark{} for CPU and GPU rendering. LA~v2 is the current
deep-zoom approximation implementation; its motivation and differences
from v1 are discussed in \cref{sec:la-v2-motivation}.

Bilinear approximation~\cite{Imagina,Thompson2023BLA} precomputes composable linear
maps that advance the perturbation delta by multiple iterations at once,
avoiding per-step evaluation when the linearization error remains small
(\cref{fig:bla-v1-hierarchy}).

\subsection{Reference orbit and delta formulation}
The perturbation framework (\cref{sec:perturbation-concept}) decomposes
the orbit of a nearby pixel \(c = c_\star + \Delta c\) into a
precomputed reference orbit \(z_n^\star\) and a delta
\(\Delta z_n = z_n - z_n^\star\).  The exact delta recurrence
(\cref{eq:perturb-exact}) is
\[
\Delta z_{n+1} = 2z_n^\star\,\Delta z_n + (\Delta z_n)^2 + \Delta c.
\]
When \(\Delta z_n\) remains small, the quadratic term can be neglected,
yielding a linearized update:
\begin{equation}
\Delta z_{n+1} \approx A_n\,\Delta z_n + B_n\,\Delta c,\qquad
A_n = 2z_n^\star,\quad B_n = 1.
\label{eq:bla-v1-one-step}
\end{equation}
BLA generalizes this into precomputed multi-step maps that \emph{jump} multiple
iterations while maintaining an explicit applicability radius. The name BLA
identifies the implemented method; its evaluated map is linear in the two
complex inputs $\Delta z$ and $\Delta c$. It contains no product
$\Delta z\Delta c$. A complex coefficient rotates and scales an input, and
adding the two coefficient products combines the incoming-state and
parameter-offset contributions.

The single-update error before rounding is exactly $(\Delta z_n)^2$.
For example, with $z_n^\star=-1/2$, $\Delta z_n=0.001$, and
$\Delta c=0.001$, the exact update is $0.000001$ while the linearized
update is zero. The relative error in that canceled result is large even
though the omitted term is small in absolute magnitude. Such cancellation
explains why an applicability test is necessary and why it does not by itself
certify the final orbit error.

\subsection{What \code{BLA<T>} stores}
The reference-dependent work can be reused only if each block has a compact
description. A \code{BLA<T>} instance stores two complex coefficients:
\[
A = A_x + iA_y,\qquad B = B_x + iB_y,
\]
plus:
\begin{itemize}
  \item \code{r2}: the squared input-radius bound $R^2$ used during lookup,
  \item \code{l}: the number of Mandelbrot iterations summarized by this step.
\end{itemize}
These objects exist to accelerate \emph{escape-time evaluation} by evolving
\(\Delta z\) cheaply for many pixels, then reconstructing \(z \approx z^\star +
\Delta z\) to perform bailout checks consistent with Mandelbrot rendering.

\subsection{Applying a step: complex multiply-add}
The method \code{getValue} (taking arguments
\code{Real\-Delta\-SubN}, \code{Imag\-Delta\-SubN},
\code{Real\-Delta\-Sub0}, \code{Imag\-Delta\-Sub0})
applies:
\[
\Delta z \leftarrow A\,\Delta z + B\,\Delta c,
\]
expanded into real arithmetic:
\begin{align}
\Re(\Delta z') &= A_x \Re(\Delta z) - A_y \Im(\Delta z) + B_x \Re(\Delta c) - B_y \Im(\Delta c), \\
\Im(\Delta z') &= A_x \Im(\Delta z) + A_y \Re(\Delta z) + B_x \Im(\Delta c) + B_y \Re(\Delta c).
\end{align}
Both expressions use the incoming real and imaginary components; neither may
use a component overwritten by the other expression. For a small complex
example, $A=i$, $B=2$, $\Delta z=1+i$, and $\Delta c=0.1-0.2i$ give
$A\Delta z=-1+i$, $B\Delta c=0.2-0.4i$, and
$\Delta z'=-0.8+0.6i$. The same arithmetic evaluates a long block once its
coefficients have been constructed.

\subsection{Composing steps to build longer jumps}
If a first step maps \(\Delta z \mapsto A_1\Delta z + B_1\Delta c\) and a second maps
\(\Delta z' \mapsto A_2\Delta z' + B_2\Delta c\), substitution gives
\begin{align*}
\Delta z'' &\approx A_2(A_1\Delta z+B_1\Delta c)+B_2\Delta c\\
            &= (A_2A_1)\Delta z+(A_2B_1+B_2)\Delta c.
\end{align*}
The resulting coefficients are
\begin{equation}
A_{\text{new}} = A_2 A_1,\qquad B_{\text{new}} = A_2 B_1 + B_2.
\label{eq:bla-v1-composition}
\end{equation}
The second block multiplies the first block's entire output. Its $A_2$ therefore
multiplies both $A_1$ and $B_1$, whereas its own parameter contribution $B_2$
is added afterward. Composition is exact for the two linear maps, but the maps
already omit nonlinear terms. If their omitted contributions are $e_1,e_2$,
the combined omission includes $A_2e_1+e_2$; a large coefficient can amplify
earlier error.
In the code (\code{BLA.h}), these two steps are passed as the parameters
\code{x} and \code{y} to \code{getNewA} and \code{getNewB}, corresponding to
step~1 and step~2 respectively.
The lengths add, $\ell_{\mathrm{new}}=\ell_1+\ell_2$.
Pairwise composition builds a binary hierarchy with nominal lengths
$1,2,4,8,\ldots$. An unpaired final entry is carried forward, so a
terminal map may cover a shorter span. The implementation constructs
the one- and two-iteration maps recursively but stores tables starting
at level~2, with nominal length~4. The retained levels therefore begin
at $4,8,16,\ldots$, as shown in \cref{fig:bla-v1-hierarchy}.

\paragraph{Worked four-iteration construction.}
\label{sec:bla-v1-worked-map}
Use $c_\star=-1/2$ and start at reference index~$r=1$. The reference
recurrence gives the following samples and one-update coefficients:
\begin{center}
\begin{tabular}{cccc}
\hline
$r$ & $z_r^\star$ & $A_r=2z_r^\star$ & $B_r$\\
\hline
1 & $-1/2$ & $-1$ & 1\\
2 & $-1/4$ & $-1/2$ & 1\\
3 & $-7/16$ & $-7/8$ & 1\\
4 & $-79/256$ & $-79/128$ & 1\\
\hline
\end{tabular}
\end{center}
The endpoint is $z_5^\star=-26527/65536$. Combining the first two
updates gives $A_{1\to3}=1/2$, $B_{1\to3}=1/2$; the next pair gives
$A_{3\to5}=553/1024$, $B_{3\to5}=49/128$. Combining those pairs yields
\begin{equation}
A_{1\to5}=\frac{553}{2048},\qquad
B_{1\to5}=\frac{553}{1024}\frac12+\frac{49}{128}
         =\frac{1337}{2048},\qquad \ell=4.
\label{eq:bla-v1-worked-coefficients}
\end{equation}
Thus four individual linearized updates reduce to
$\Delta z_{\mathrm{out}}\approx(553\Delta z+1337\Delta c)/2048$.
The block's reconstructed endpoint is
$z_5^\star+\Delta z_{\mathrm{out}}$, not the reference sample at its input.

For a fresh neighboring pixel, the exact first update gives
$\Delta z_1=\Delta c=t$. With $t=0.001$, the four-update BLA estimate
is $0.0009228515625$, while repeated exact perturbation gives
approximately $0.000921980959229$. The difference is approximately
$8.71\cdot10^{-7}$. The deliberately large offset makes the omission visible;
it is not an example of an accepted production BLA block. The radius calculation
below shows how a much smaller offset can pass lookup. The same reference and
four-update span are used for the LA~v2 comparison in
\cref{sec:la-worked-shared-map}.

\subsection{Constructing and composing validity radii}
\label{sec:bla-v1-validity}

The one-iteration map drops $(\Delta z)^2$. To keep that term small
relative to $A_n\Delta z$, the constructor assigns
\begin{equation}
R_n=\epsilon |A_n|=2\epsilon |z_n^\star|,
\qquad \epsilon=2^{-23},\qquad \code{r2}=R_n^2.
\label{eq:bla-v1-radius}
\end{equation}
The value of $\epsilon$ comes from \code{BLAS::BLA\_BITS}; the same
construction setting is used for the supported scalar types. The
radius is an applicability threshold for the omitted quadratic term,
rather than a certified bound on the complete pixel orbit error.
For nonzero $A_n\Delta z$, the ratio of the omitted quadratic magnitude to
the linear magnitude is $|\Delta z|/|A_n|$. Requiring
$|\Delta z|<\epsilon|A_n|$ keeps that ratio below~$\epsilon$. The parameter
term can cancel the linear term, and errors from previous blocks can accumulate,
so that local comparison is not a relative-error bound on the complete result.

For a first map with coefficients $A_1,B_1$ and radius $R_1$, followed
by a second map with radius $R_2$, the implementation uses the stored
maximum parameter-offset radius $R_{\max}$ to compose the bounds.
For $|A_1|>0$, the merged radius is
\begin{equation}
R_{\mathrm{new}}=\min\!\left(
  R_1,\ \max\!\left(0,\frac{R_2-|B_1|R_{\max}}{|A_1|}\right)
\right).
\label{eq:bla-v1-merged-radius}
\end{equation}
The first bound keeps the incoming delta within the first map's domain.
The second follows from
$|A_1\Delta z+B_1\Delta c|\le |A_1||\Delta z|+|B_1|R_{\max}$
and keeps the intermediate delta within the second domain. If the
parameter contribution consumes that domain, the merged radius becomes
zero and the strict lookup test rejects the map.

\paragraph{Worked radius composition.}
Use the coefficients from \cref{sec:bla-v1-worked-map} and choose a view
whose maximum parameter offset is $R_{\max}=\epsilon/8$. Expressing every
radius in units of~$\epsilon$ keeps the calculation readable. The one-update
radii are $1,1/2,7/8,79/128$ in those units. For the first pair,
\[
\frac{R_{1\to3}}{\epsilon}
 =\min\!\left(1,\frac{1/2-1/8}{1}\right)=\frac38.
\]
For the second pair,
\[
\frac{R_{3\to5}}{\epsilon}
 =\min\!\left(\frac78,\frac{79/128-1/8}{7/8}\right)=\frac9{16}.
\]
Finally, the first pair has $|A|=|B|=1/2$, so
\[
\frac{R_{1\to5}}{\epsilon}
 =\min\!\left(\frac38,
     \frac{9/16-(1/2)(1/8)}{1/2}\right)=\frac38.
\]
An incoming $\Delta z=\Delta c=\epsilon/16$ therefore satisfies the
strict four-update lookup test. Increasing the view radius consumes more of
the intermediate domain even if the incoming delta is unchanged. At a zero
reference, the one-update radius from \eqref{eq:bla-v1-radius} is zero;
the strict test cannot accept that one-update map. Formula
\eqref{eq:bla-v1-merged-radius} is written for a nonzero denominator, rather
than defining division by a zero coefficient.

\subsection{Aligned lookup and perturbation fallback}
\label{sec:bla-v1-lookup}

The CPU \code{BLAS} and GPU \code{GPU\_BLAS} implementations select a
map using the current reference index $r$ and squared delta magnitude
$|\Delta z|^2$. The accumulated pixel count $n$ is separate from $r$,
because rebasing can change the reference position while preserving
completed work. Tables begin at reference index~1; a nominal
$2^k$-iteration map is aligned when $r-1$ is divisible by $2^k$.
The longest candidate level is also limited by the available table.
Lookup scans from that level toward the stored four-iteration level
and accepts the first map satisfying
\begin{equation}
|\Delta z|^2<\code{r2}.
\label{eq:bla-v1-lookup}
\end{equation}
Equality is rejected. If no aligned retained map passes, the renderer
performs an individual update using the exact perturbation recurrence
and can attempt lookup again at the new reference position.

\begin{figure}[!htbp]
\centering
\begin{tikzpicture}[>=Stealth, x=0.58cm, y=1cm,
    every node/.style={font=\small},
    block/.style={draw=gray!65, fill=gray!8},
    rejected/.style={draw=red!70!black, fill=red!6, dashed},
    accepted/.style={draw=green!45!black, fill=green!10},
    perturbation/.style={->, thick, orange!80!black}]
  % Table coordinates use r-1 so all retained blocks share the actual alignment.
  \draw[->, gray] (0,0) -- (16.5,0);
  \foreach \x/\r in {0/1,4/5,8/9,12/13,16/17} {
    \draw (\x,-0.08) -- (\x,0.08);
    \node[below] at (\x,-0.08) {$\r$};
  }
  \node[anchor=east] at (-0.3,0) {Reference $r$};
  \node[anchor=east] at (-0.3,0.85) {$\ell=4$};
  \node[anchor=east] at (-0.3,1.75) {$\ell=8$};
  \node[anchor=east] at (-0.3,2.65) {$\ell=16$};
  \draw[block] (0,0.55) rectangle (4,1.15);
  \draw[block] (4,0.55) rectangle (8,1.15);
  \draw[rejected] (8,0.55) rectangle (12,1.15);
  \draw[accepted] (12,0.55) rectangle (16,1.15);
  \draw[accepted] (0,1.45) rectangle (8,2.05);
  \node at (4,1.75) {Selected at $r=1$};
  \draw[rejected] (8,1.45) rectangle (16,2.05);
  \node at (12,1.75) {Invalid at $r=9$};
  \draw[rejected] (0,2.35) rectangle (16,2.95);
  \node at (8,2.65) {Larger candidate invalid at $r=1$};
  \node[anchor=east] at (-0.3,-1.0) {Traversal};
  \draw[->, thick, green!45!black] (0,-1.0) -- (8,-1.0);
  \foreach \a/\b in {8/9,9/10,10/11,11/12} {
    \draw[perturbation] (\a,-1.0) -- (\b,-1.0);
  }
  \draw[->, thick, green!45!black] (12,-1.0) -- (16,-1.0);
  \node[below, font=\scriptsize] at (4,-1.15) {$\ell=8$};
  \node[below, font=\scriptsize, align=center] at (10,-1.15)
    {Four individual\\updates};
  \node[below, font=\scriptsize] at (14,-1.15) {$\ell=4$};
\end{tikzpicture}
\caption[BLA v1 hierarchy and fallback]{Illustrative BLA~v1 lookup and
  traversal in a sufficiently long reference. At $r=1$, the
  sixteen-iteration candidate fails the radius test and the
  eight-iteration candidate passes. At $r=9$, both retained aligned
  candidates fail; four individual updates reach the next
  four-iteration boundary at $r=13$, where a map passes. Colors describe
  validity at each encountered pixel state, not an intrinsic property
  of a table entry. The reference and pixel budgets exceed the displayed
  endpoint, and no bailout or rebasing occurs in this schematic.}
\label{fig:bla-v1-hierarchy}
\end{figure}

Radius validity alone does not authorize a jump. The CPU and CUDA
renderers separately require the endpoint to exist in the stored
reference and reject a proposed map when $n+\ell\ge n_{\max}$.
A map that exactly reaches the pixel limit therefore also falls back
to individual updates. These endpoint and budget checks occur after
lookup: rejection leaves the delta and counts unchanged, without
requesting a smaller map for that lookup. LA~v2 instead descends to a
finer stage when a block does not fit, and permits an exact budget fit
(\cref{sec:la-stage-hierarchy}).

An accepted BLA map advances both $n$ and $r$ by $\ell$, evaluates
$A\Delta z+B\Delta c$, and updates the squared delta magnitude for
the next lookup. Rendering reconstructs the absolute state at the
points needed for bailout and rebasing; the exact placement of those
tests differs between the native-double and HDR CUDA kernels. When
rebasing is required, the absolute orbit estimate becomes the delta
relative to the reference origin while $n$ is preserved.

\paragraph{Worked count and reference-index trace.}
Consider the validity pattern in \cref{fig:bla-v1-hierarchy}, a reference long
enough to contain every shown endpoint, and $n_{\max}=21$. Suppose no bailout or
rebase occurs during these updates. After the initial individual update,
$(n,r)=(1,1)$. At $r=17$, assume that longer candidates are unavailable or fail
their radius tests while the four-update candidate passes its radius test.
The lookup therefore selects that candidate before checking its budget. The
lookup and count changes are
\begin{center}
\begin{tabular}{clc}
\hline
Incoming $(n,r)$ & Action & Outgoing $(n,r)$\\
\hline
$(1,1)$ & Length 16 fails radius; length 8 passes & $(9,9)$\\
$(9,9)$ & Candidates fail; four individual updates & $(13,13)$\\
$(13,13)$ & Length 4 passes radius and budget & $(17,17)$\\
$(17,17)$ & Length 4 reaches 21 exactly; reject & $(17,17)$\\
$(17,17)$ & Four individual updates finish the budget & $(21,21)$\\
\hline
\end{tabular}
\end{center}
The last rejection illustrates the strict BLA budget rule rather than a radius
failure. No coefficients are applied before the individual tail begins. For
comparison, an LA~v2 block that reaches the budget exactly can be accepted.
If rebasing instead occurs at $(n,r)=(17,17)$, reconstruction first forms
$z=z_{17}^\star+\Delta z$; assigning $\Delta z=z$ and $r=0$ preserves $n=17$.
The next individual update reaches $(n,r)=(18,1)$ if no bailout occurs.
Both counts advance during a block, while only the reference count resets
during rebasing.

\subsection{Construction, storage, and evaluation costs}
The hierarchy trades preprocessing and reference-dependent storage for fewer
per-pixel recurrence evaluations. Each retained entry contains four real
coefficient components, one real squared radius, and one integer length.
Nominal four-update entries are approximately $n_{\mathrm{ref}}/4$ in number,
eight-update entries approximately $n_{\mathrm{ref}}/8$, and so on. Ignoring
terminal rounding, their geometric sum is less than
$n_{\mathrm{ref}}/2$. Construction therefore requires linear work and storage
in the reference length, with constants determined by arithmetic type and
memory layout.

Each accepted block costs two complex products and an addition, plus lookup
and count management, regardless of its span. The saved work increases with
accepted length. Long tables do not guarantee long accepted jumps: alignment,
radius, endpoint, and budget restrictions can all force individual updates.
GPU performance also depends on table traffic and how consistently neighboring
threads choose the same path. Those costs explain why both the reference and
the pixel domain matter when comparing acceleration methods.

\subsection{Limits that motivate LA~v2}

Near a return to the critical point, $|z_r^\star|$ can be very small.
\eqref{eq:bla-v1-radius} then gives a restrictive radius, and a composed
block crossing that return can inherit a restrictive intermediate
bound. A fixed binary hierarchy cannot move its boundaries to place
the return at a block's nonlinear input. BLA~v1 consequently performs
individual updates wherever no retained aligned map is applicable.

LA~v2 retains the quadratic term at each block input and chooses useful
boundaries from the orbit's dynamics. The resulting staged maps can
remain useful over longer spans at deep zooms, as explained in
\cref{sec:la-v2-motivation}. BLA~v1 remains available for comparison
and can perform well at shallow zooms, where elaborate approximation
tables may provide little benefit. Neither method's hierarchy alone
establishes a universal performance or accuracy ordering.
